\chapter*{ประวัติคณะผู้เขียน}
\addcontentsline{toc}{chapter}{ประวัติคณะผู้เขียน}

\begin{center}
\begin{minipage}[c]{0.26\textwidth}
    \centering
    \begin{tikzpicture}
        % กรอบวงกลมภาพถ่ายผู้เขียนขลิบทองราชภัฏ
        \draw[color=rbruGold, line width=2.5pt, fill=rbruNavy!15] (0,0) circle (1.8cm);
        \node at (0,0) {\fontsize{11}{14}\selectfont\bfseries\color{rbruNavy} ภาพถ่ายผู้เขียน};
    \end{tikzpicture}
\end{minipage}
\hfill
\begin{minipage}[c]{0.70\textwidth}
    {\fontsize{18}{22}\selectfont\bfseries\color{rbruNavy} ผู้ช่วยศาสตราจารย์ ดร.ชีวะ ทัศนา\par}
    \vspace{2pt}
    {\fontsize{12}{15}\selectfont\bfseries\color{durianGold} ผู้นิพนธ์หลักและผู้อำนวยการหลักสูตร AIoT\par}
    \vspace{4pt}
    {\fontsize{11}{14}\selectfont สาขาวิชาฟิสิกส์และวิทยาศาสตร์ทั่วไป คณะวิทยาศาสตร์และเทคโนโลยี\par}
    {\fontsize{11}{14}\selectfont มหาวิทยาลัยราชภัฏรำไพพรรณี จันทบุรี\par}
    {\fontsize{10}{13}\selectfont อีเมล: \texttt{chewa.t@rbru.ac.th} \textbar\ Google Scholar: \texttt{jJz3B6AAAAAJ}\par}
\end{minipage}
\end{center}

\vspace{0.4cm}
{\color{rbruGold}\hrule height 1.2pt}
\vspace{0.4cm}

\begin{minipage}[t]{0.48\textwidth}
    {\fontsize{13}{16}\selectfont\bfseries\color{rbruNavy} ประวัติการศึกษา}\\[6pt]
    \begin{itemize}[leftmargin=1.2em, itemsep=2pt, topsep=0pt]
        \item \textbf{ปร.ด. (ฟิสิกส์ประยุกต์)} สถาบันเทคโนโลยีพระจอมเกล้าเจ้าคุณทหารลาดกระบัง (KMITL)
        \item \textbf{วท.ม. (ฟิสิกส์ประยุกต์)} สถาบันเทคโนโลยีพระจอมเกล้าเจ้าคุณทหารลาดกระบัง (KMITL)
        \item \textbf{วท.บ. ฟิสิกส์ (คอมพิวเตอร์และอิเล็กทรอนิกส์)} มหาวิทยาลัยนเรศวร
    \end{itemize}

    \vspace{0.3cm}
    {\fontsize{13}{16}\selectfont\bfseries\color{rbruNavy} ความเชี่ยวชาญทางวิชาการ}\\[6pt]
    \begin{itemize}[leftmargin=1.2em, itemsep=2pt, topsep=0pt]
        \item ฟิสิกส์เกษตรดิจิทัลและเคมีไฟฟ้าของดิน
        \item สมองกลฝังตัวประสิทธิภาพสูง (ESP32-S3)
        \item ปัญญาประดิษฐ์ฝังตัวที่ขอบ (Edge TinyML)
        \item สวนศาสตร์การสั่นสะเทือนและการวิเคราะห์ FFT
    \end{itemize}
\end{minipage}
\hfill
\begin{minipage}[t]{0.48\textwidth}
    {\fontsize{13}{16}\selectfont\bfseries\color{rbruNavy} คณะวิทยากรผู้ร่วมจัดทำตำรา}\\[6pt]
    \begin{itemize}[leftmargin=1.2em, itemsep=2pt, topsep=0pt]
        \item \textbf{ผศ.อรรถกร คำฉัตร} วท.ม. (เทคโนโลยีสารสนเทศ) เชี่ยวชาญ โครงข่าย LoRaWAN คลาวด์ IoT และระบบอัตโนมัติ
        \item \textbf{ผศ.ดร.ณมนรัก คำฉัตร} ปร.ด. (ชีววิทยา/พืชศาสตร์) เชี่ยวชาญ สรีรวิทยาพืชและการประเมิน $ET_c$ ทุเรียน
        \item \textbf{รศ.ดร.นิภัทร เปี่ยมอรุณ} ปร.ด. (เคมีประยุกต์) เชี่ยวชาญ เคมีเชิงฟิสิกส์และการวิเคราะห์สเปซสี HSV
        \item \textbf{ผศ.ดร.จิรภัทร จันทมาลี} ปร.ด. (จุลชีววิทยา) เชี่ยวชาญ โรคพืชทุเรียนและจุลินทรีย์ก่อโรคในสัตว์น้ำ
        \item \textbf{ผศ.นิคม รัตนโรจนกุล} วท.ม. (ฟิสิกส์ประยุกต์) เชี่ยวชาญ ระบบพลังงานแสงอาทิตย์และแบตเตอรี่ LiFePO4
        \item \textbf{รศ.ดร.นันทพร มูลรังษี} ปร.ด. (เคมีวิเคราะห์) เชี่ยวชาญ สารพฤกษเคมีและมาตรฐานคุณภาพผลผลิต
        \item \textbf{ผศ.ดร.สุพัตรา รักษาพรต} ปร.ด. (เทคโนโลยีเคมี) เชี่ยวชาญ การสอบเทียบเซนเซอร์เคมีและคุณภาพน้ำ
    \end{itemize}
\end{minipage}

\vfill
\begin{center}
    {\fontsize{10}{13}\selectfont\color{slateGrey} คณะวิทยาศาสตร์และเทคโนโลยี มหาวิทยาลัยราชภัฏรำไพพรรณี เลขที่ 41 หมู่ 5 ต.ท่าช้าง อ.เมือง จ.จันทบุรี 22000}
\end{center}
\clearpage
